\documentclass[11pt]{article}
\usepackage[a4paper,margin=30mm]{geometry}
\usepackage{amsmath,amssymb,amsthm,mathtools,array}
\usepackage[hidelinks]{hyperref}
\usepackage{xcolor}

\newtheorem{theorem}{Theorem}
\newtheorem{lemma}[theorem]{Lemma}
\newtheorem{proposition}[theorem]{Proposition}
\newtheorem{corollary}[theorem]{Corollary}
\theoremstyle{definition}
\newtheorem{definition}[theorem]{Definition}
\newtheorem{hypothesis}[theorem]{External interface}
\theoremstyle{remark}
\newtheorem{remark}[theorem]{Remark}

\newcommand{\R}{\mathbb R}
\newcommand{\T}{\mathbb R/\mathbb Z}
\newcommand{\mesh}{\operatorname{mesh}}
\newcommand{\Ind}{\operatorname{Ind}}
\newcommand{\diam}{\operatorname{diam}}
\newcommand{\status}[1]{\par\smallskip\noindent\textbf{Status:}
  \begingroup\small\sffamily #1.\endgroup\par\smallskip}

\title{Inscribed Squares for Jordan Curves of Vanishing Critical
  $2$-Variation}
\author{Clanker contributors\\
\small Issue \#134; verification-critical article for PR \#135}
\date{30 August 2026}

\begin{document}
\maketitle

\begin{abstract}
We prove that a Jordan parametrization satisfying the critical vanishing
$2$-variation condition $V_0$ inscribes a positive-size square with four
distinct vertices.  The proof combines finite-partition approximation,
anchored regular smoothing, an excursion-current winding estimate, and
compactness of Liouville primitives.  Its final implication is conditional
on the precisely stated and version-pinned Boedihardjo--Geng and
Asano--Ike interfaces, together with the cited index and Green identities.
This is not a proof of the unrestricted Square Peg conjecture and is not an
end-to-end \texttt{verified:lean} result.
\end{abstract}

\begin{center}
\fbox{\parbox{0.88\textwidth}{\centering
\textbf{VERIFIED CONDITIONAL ON CITED INTERFACES}.  This audit verdict is
neither a grant of \texttt{verified:review} nor an end-to-end
\texttt{verified:lean} claim.}}
\end{center}

\section{Main result and notation}

\begin{definition}[Jordan parametrization]
A Jordan parametrization is a continuous, $1$-periodic map
$c\colon\R\to\R^2$ whose restriction to $[0,1)$ is injective.  Equivalently,
it is an embedding $\T\to\R^2$ together with its periodic lift.  Four
vertices are distinct when they are pairwise distinct image points, and a
square has positive size when its side length is positive.
\end{definition}

\begin{definition}[Partitions and variation]
For an interval $I=[a,b]$, an ordered partition is a finite sequence
$a=t_0<\cdots<t_m=b$.  For such partitions define
\[
 \|x\|_{2\text{-var};I}^2
 :=\sup_{a=t_0<\cdots<t_m=b}
       \sum_{j=0}^{m-1}|x(t_{j+1})-x(t_j)|^2.
\]
For a partition $D=\{0=t_0<\cdots<t_m=1\}$, put
\[
 \mesh(D)=\max_j(t_{j+1}-t_j),\qquad
w_2(c,\delta)^2
 :=\sup_{\mesh(D)\leq\delta}
       \sum_j|c(t_{j+1})-c(t_j)|^2.
\]
Both suprema initially take values in $[0,\infty]$.
Thus $w_2(c,\delta)$ is the nonnegative extended square root of the displayed
supremum.  It is used as a real number only in the finite regime supplied by
$V_0$ below.
\end{definition}

\begin{definition}[Critical vanishing-variation hypothesis]
Write $V_0(c)$ for
\[
 \exists\delta_0>0:\ w_2(c,\delta_0)<\infty
 \quad\text{and}\quad w_2(c,\delta)\longrightarrow0
                 \quad(\delta\downarrow0).
\]
Thus every real-valued variation calculation below occurs in the finite
regime $0<\delta\leq\delta_0$.
\end{definition}

In this document the repository shorthand $C^{0,2\text{-var}}$ means exactly
the critical vanishing condition $V_0$ above.  The theorem is stated only
under $V_0$ and makes no separate claim that the global $2$-variation is
finite.

\begin{theorem}[Main theorem]\label{thm:main}
Let $c\colon\R\to\R^2$ be a $1$-periodic Jordan parametrization satisfying
$V_0(c)$.  Then four pairwise distinct points of $c(\R)$ are the vertices of a
positive-size Euclidean square.
\status{PROVED from the internal reductions and the independently checked
cited interfaces; conditional on those interfaces}
\end{theorem}

\begin{proof}
This is Corollary~\ref{cor:positive-square}, proved after the
approximation and primitive-compactness hypotheses of the cited
Asano--Ike criterion have been discharged.
\end{proof}

\section{Polygonal approximation}

\begin{lemma}[Lifted short-interval modulus]\label{lem:local-modulus}
If $V_0(c)$ holds, $0<\delta\leq\tfrac12$, and $I\subset\R$ has
$|I|\leq\delta$, then
\[
       \|c\|_{2\text{-var};I}\leq\sqrt2\,w_2(c,\delta).
\]
\status{PROVED}
\end{lemma}

\begin{proof}
If $I$ does not cross an integer, extend any partition of $I$ to a
full-period partition of mesh at most $\delta$.  If it crosses one integer,
insert the seam.  The insertion replaces at most one increment $u+v$ by
$u,v$, and
\(
 |u+v|^2\leq2(|u|^2+|v|^2).
\)
After reducing the two pieces modulo one, fill every complementary gap by
subintervals of length at most $\delta$ and combine the nodes into a
full-period partition.  Since $|I|\leq\delta\leq1/2$, at most one integer
seam occurs.  Take suprema.
\end{proof}

For a partition $D$, let $c^D$ be the affine interpolant of $c$ on every
interval of $D$.

\begin{lemma}[Every-partition interpolation estimate]
For every partition $D$ of $[0,1]$,
\[
 \|c-c^D\|_{2\text{-var};[0,1]}
 \leq \sqrt8\,w_2(c,\mesh(D)).                         \tag{2.1}
\]
Consequently $V_0(c)$ implies $c^D\to c$ in $2$-variation along
\emph{every} sequence with $\mesh(D)\to0$.
\status{PROVED; replaces the load-bearing Yang quantifier}
\end{lemma}

\begin{proof}
Write $I_j=[t_j,t_{j+1}]$ and $y=c-c^D$.  Since $y$ vanishes at every
$t_j$, inserting all partition nodes into an arbitrary variation partition
costs at most a factor $2$ in the squared sum.  Indeed, although one
crossing increment may meet many nodes of $D$, $y=0$ at every inserted
node, so only its two endpoint pieces can be nonzero.  Hence
\[
 \|y\|_{2\text{-var}}^2
 \leq2\sum_j\|y\|_{2\text{-var};I_j}^2.
\]
On $I_j$, Minkowski's inequality and the fact that the affine chord has
$2$-variation equal to $|c(t_{j+1})-c(t_j)|$ give
\[
 \|y\|_{2\text{-var};I_j}
 \leq2\|c\|_{2\text{-var};I_j}.
\]
Near-optimal local partitions concatenate into one partition of mesh at
most $\mesh(D)$, so
\[
 \sum_j\|c\|_{2\text{-var};I_j}^2
 \leq w_2(c,\mesh(D))^2.
\]
Combining the three inequalities proves (2.1).
\end{proof}

\begin{hypothesis}[Simple interpolation, BG]\label{hyp:bg}
For every $\varepsilon>0$ and every Jordan parametrization $c$, there is a
partition $D$ with $\mesh(D)<\varepsilon$ such that the parametrized affine
interpolant $c^D$ is a simple polygonal Jordan parametrization.
\status{CITED: Boedihardjo--Geng~\cite{BG}, arXiv:1309.1576v2, Theorem~2.2,
PDF pp.~6--7; the fixed-parameter, parameter-mesh, unique-minimizing-geodesic,
and Jordan-interpolant clauses were independently checked.  In $\R^2$ the
geodesics are affine chords}
\end{hypothesis}

Under Interface~\ref{hyp:bg} and $V_0(c)$, fix
$\delta_0>0$ with $w_2(c,\delta_0)<\infty$ as furnished by the definition,
then choose $D_n$ with
$\mesh(D_n)<\min(\delta_0,1/n)$ and put $P_n=c^{D_n}$.  The preceding proved
lemma gives
\[
 \|P_n-c\|_{2\text{-var}}\to0,
 \qquad P_n\to c\text{ uniformly},                       \tag{2.2}
\]
the uniform statement following because $P_n(0)=c(0)$.

\section{Regular parametrized smoothing}

\begin{lemma}[Regular parametrized corner smoothing]\label{lem:smooth}
For every parametrized simple polygon $P\colon\T\to\R^2$ and every
$\varepsilon>0$, there is a $C^\infty$ \emph{regular} Jordan embedding
$Q\colon\T\to\R^2$ such that
\[
 Q(0)=P(0),\qquad
 \|Q-P\|_{1\text{-var}}<\varepsilon.                   \tag{3.1}
\]
Here regular means $Q'(t)\ne0$ at every parameter, including the periodic
seam.
\status{PROVED}
\end{lemma}

\begin{proof}
Let $0=t_0<\cdots<t_m=1$ be the edge partition on a periodic lift, let
$v_i$ be the nonzero constant velocity on the $i$th edge, and put
$M=\max_i|v_i|$.  Choose an even, nonnegative, smooth periodic mollifier
$\rho_h$, supported at circular distance less than $h$ and having integral
one, where $2h$ is smaller than every edge's parameter length.  Set
$Q_h=P*\rho_h$.  Writing
\(
Q_h(t)=\int_{\T}P(r)\rho_h(t-r)\,dr,
\)
differentiation of the smooth kernel under the integral makes $Q_h$ smooth.
Periodic integration by parts, using absolute continuity of $P$, gives
\[
 Q_h'=P'*\rho_h.                                       \tag{3.2}
\]
Thus $Q_h$ is periodic and smooth, including at $0=1$.  Away from the
$h$-neighborhoods of vertices, $Q_h'=P'=v_i$.  Near the vertex between
$v_{i-1}$ and $v_i$, equation (3.2) is a convex combination
$\alpha v_{i-1}+\beta v_i$, with $\alpha,\beta\geq0$ and
$\alpha+\beta=1$.

The adjacent velocities cannot lie on opposite rays, since then the two
incident polygonal edges would overlap near their common vertex.  Hence a
linear functional $\ell_i$ is strictly positive on both.  It follows that
$\ell_i(Q_h')>0$ throughout the vertex neighborhood.  Therefore $Q_h'$ never
vanishes.

To keep the injectivity quantifiers uniform, first fix
$h_0<\tfrac14\min_i(t_{i+1}-t_i)$ and a finite cover of $\T$ by open vertex
parameter intervals of radius $2h_0$ and open intervening edge-interior
parameter intervals.  Shrink
$h_0$ once, if necessary, so the positivity argument applies on each fixed
vertex interval for every $0<h<h_0$; on the remaining pieces $Q_h'$ is an
edge velocity.  Hence $Q_h$ is injective on every member of this one fixed
cover for every $0<h<h_0$.

Let $\eta>0$ be a Lebesgue number for the fixed cover.  On the compact set of
parameter pairs at circular distance at least $\eta$, injectivity of $P$
gives a positive minimum separation $d_\eta$.  Moreover,
\[
 |Q_h(t)-P(t)|\leq
 \sup_{d_{\T}(s,0)<h}|P(t-s)-P(t)|,
\]
which tends uniformly to zero.  Choose $h$ so that
$\|Q_h-P\|_\infty<d_\eta/3$.  Such pairs remain separated by $Q_h$; pairs
closer than $\eta$ lie in a common cover member and are separated by local
injectivity.  Hence $Q_h$ is a regular Jordan embedding.

Finally, $Q_h-P$ is absolutely continuous, so its $1$-variation is the
integral of $|Q_h'-P'|$.  This integrand vanishes off the union $U_h$ of the
$m$ vertex neighborhoods, while $|U_h|\leq2mh$ and both derivatives have
norm at most $M$ there.  Consequently
\[
 \|Q_h-P\|_{1\text{-var}}
 =\int_0^1|Q_h'-P'|\leq4mMh.
\]
If orientation is part of
the smooth-Jordan convention, it is preserved: choose a point in the bounded
component of $\R^2\setminus P(\T)$ and
then $h$ so small that the straight homotopy from $P$ to $Q_h$ avoids it;
winding-number homotopy invariance fixes the sign.

Now translate the mollification by setting
\[
 \widetilde Q_h(t)=Q_h(t)+P(0)-Q_h(0).
\]
It is periodic and $C^\infty$, and $\widetilde Q_h'=Q_h'$, so it is regular.
Translation preserves injectivity, the Jordan property, and orientation, and
$\widetilde Q_h(0)=P(0)$.  Moreover
$(\widetilde Q_h-P)'=(Q_h-P)'$, whence
\[
 \|\widetilde Q_h-P\|_{1\text{-var}}
 =\|Q_h-P\|_{1\text{-var}}\leq4mMh.                    \tag{3.3}
\]
Because $\widetilde Q_h-P$ vanishes at $0$, absolute continuity also gives,
for every $t\in[0,1]$,
\[
 |\widetilde Q_h(t)-P(t)|
 \leq\int_0^t|\widetilde Q_h'-P'|
 \leq\|\widetilde Q_h-P\|_{1\text{-var}}.             \tag{3.4}
\]
Thus $\|\widetilde Q_h-P\|_\infty\leq
\|\widetilde Q_h-P\|_{1\text{-var}}$, and taking also
$h<\varepsilon/(4mM)$ proves the lemma with $Q=\widetilde Q_h$.
\end{proof}

Using Lemma~\ref{lem:smooth}, choose $c_n$ from $P_n$ with errors
$\varepsilon_n\downarrow0$, anchored by $c_n(0)=P_n(0)=c(0)$.  Since
$\|x\|_{2\text{-var}}\leq\|x\|_{1\text{-var}}$, (2.2) gives
\[
 \|c_n-c\|_{2\text{-var}}\to0.                         \tag{3.5}
\]
The anchoring and the same variation estimate also give
\[
 \|c_n-P_n\|_\infty
 \leq\|c_n-P_n\|_{1\text{-var}}<\varepsilon_n,
\]
so (2.2) and the triangle inequality show that $c_n\to c$ uniformly.

\section{Uniform local variation}

Define, for $0<\delta\leq\tfrac12$,
\[
 \nu(\delta)=
 \sup_n\sup_{0\leq t-s\leq\delta}
       \|c_n\|_{2\text{-var};[s,t]},                    \tag{4.1}
\]
where intervals are interpreted on the lifted periodic parameter.

\begin{lemma}[Whole-family local modulus]\label{lem:nu}
Under (3.5), $\nu(\delta)\to0$ as $\delta\downarrow0$.
\status{PROVED}
\end{lemma}

\begin{proof}
Put $e_n=\|c_n-c\|_{2\text{-var};[0,1]}$.  Restriction, the seam insertion
argument, and Lemma~\ref{lem:local-modulus} give
\[
 \|c_n\|_{2\text{-var};[s,t]}
 \leq\sqrt2\bigl(w_2(c,t-s)+e_n\bigr).                  \tag{4.2}
\]
Given $\eta>0$, first choose $N$ so the right side is small for $n\geq N$,
specifically $\sqrt2e_n<\eta/2$, then choose $\delta$ using $V_0(c)$ so
$\sqrt2w_2(c,\delta)<\eta/2$.  For the finite prefix $n<N$, each
regular $C^1$ map is Lipschitz, say with constant $L_n$, and
\[
 \sum_i|\Delta c_n|^2
 \leq(\max_i|\Delta c_n|)\sum_i|\Delta c_n|
 \leq L_n^2\delta^2.
\]
Taking the finite maximum proves the claim.
\end{proof}

\section{Winding control and primitive compactness}

\begin{hypothesis}[Index-current identity]\label{hyp:index-current}
For every closed rectifiable planar curve $L$, the winding function is
compactly supported and belongs to $L^1(\R^2)$, and the current with density
$\Ind(L,\cdot)$ has boundary equal to the integration current $[L]$.
\status{CITED: integrability from Carmona--Cuf\'i~\cite{CC}, Theorem~2; boundary
identity from the main theorem of Cuf\'i--Verdera~\cite{CV} applied to smooth compactly
supported test functions.  Exact pins: J.~J. Carmona--J. Cuf\'i,
J. Analyse Math. 120 (2013), 225--253, DOI
10.1007/s11854-013-0019-9, preprint Theorem~2, PDF p.~8; and
arXiv:1306.6832v1, main theorem, PDF p.~1.  Both were independently checked}
\end{hypothesis}

Compact support here is elementary: the trace of $L$ lies in a ball, and
winding number is zero on the unbounded exterior component.  Carmona--Cuf\'i
gives $L^2$; finite measure of the supporting ball then gives $L^1$.

\begin{proposition}[Embedded-arc winding estimate]\label{prop:arc}
Using the cited Interface~\ref{hyp:index-current}, let
$a\colon[u,v]\to\R^2$ be an injective rectifiable arc.  Close it by the
oriented chord from $a(v)$ to $a(u)$, obtaining $L$.  Then
\[
 \int_{\R^2}|\Ind(L,z)|\,dz
 \leq\frac\pi4\|a\|_{2\text{-var};[u,v]}^2.             \tag{5.1}
\]
\status{PROVED from the independently checked cited index-current, Jordan,
and isodiametric theorems}
\end{proposition}

\begin{proof}
Let $H$ be its supporting line and $S$ the closing chord.  The relatively
open set
$\{t\in(u,v):a(t)\notin S\}$ is a countable disjoint union of intervals
$(\alpha_j,\beta_j)$.  Continuity puts both endpoint images on $S$.
Injectivity makes them distinct, and the excursion interior misses $S$;
hence the excursion closed by its subchord is a Jordan loop $E_j$.
Gilman--Kra--Rodr\'iguez~\cite{GKR}, Theorem~5.23, supplies its bounded complementary
domain $\Omega_j$ and
$\Ind(E_j,\cdot)=\sigma_j\mathbf1_{\Omega_j}$ almost everywhere, with
$\sigma_j\in\{\pm1\}$.

Adding the subchord does not change the excursion diameter: if $x$ is on the
subchord and $y$ is in the excursion closure, convexity of distance along
the chord gives $|x-y|\leq\max(|a(\alpha_j)-y|,|a(\beta_j)-y|)$.  Thus
$\diam(E_j)=\diam(a([\alpha_j,\beta_j]))$.  For a finite set of excursions,
choose within each parameter interval a pair arbitrarily close to realizing
that diameter.  All pairs form one ordered partition of $[u,v]$; discarding
the nonnegative increments between pairs and then taking finite and
approximation limits gives
\[
 \sum_j\diam(E_j)^2
 \leq\|a\|_{2\text{-var};[u,v]}^2.                       \tag{5.2}
\]
The planar isodiametric inequality gives
\(
 |\Omega_j|\leq(\pi/4)\diam(E_j)^2.
\)
Rectifiability and disjointness of the excursion intervals give
\[
 \sum_j\operatorname{length}(E_j)
 \leq\sum_j\bigl(\operatorname{length}(a|_{[\alpha_j,\beta_j]})
       +|a(\beta_j)-a(\alpha_j)|\bigr)
 \leq2\operatorname{length}(a),
\]
so $\sum_j[E_j]$ converges absolutely in current mass.
The difference $R=[L]-\sum_j[E_j]$ is a cycle supported on the chord line.
The boundary is zero because boundary commutes with a mass-convergent sum
when tested on differentials.
Such a cycle vanishes.  Indeed, write the chord line as $H=p_0+\R e$, where
$e$ is a unit vector.  For a compactly supported smooth one-form $\omega$,
put $g(s)=\omega_{p_0+se}(e)$ and
$G_0(s)=\int_{-\infty}^s g(r)\,dr$.  Since $R$ has compact support, choose a
longitudinal cutoff equal to one on a neighborhood of the projection of
$\operatorname{spt}R$, and a transverse cutoff equal to one near $H$ over
that projection.  Multiplying $G_0$ by these cutoffs produces
$G\in C_c^\infty(\R^2)$ with $dG(e)=g$ on
$\operatorname{spt}R\cap H$.  In the rectifiable representation of $R$, its
approximate tangent is parallel to $e$ almost everywhere because $R$ is
supported on $H$.  Hence $R(\omega-dG)=0$, while
$R(dG)=\partial R(G)=0$.  Thus $R(\omega)=0$ for every test one-form, so
$R=0$.

By (5.2) and Bieberbach's planar isodiametric inequality~\cite{Bieberbach}
(Jahresbericht DMV 24 (1915), 247--250),
$W=\sum_j\sigma_j[\Omega_j]$ converges in mass and
$\partial W=[L]$: Interface~\ref{hyp:index-current}, applied first to each
$E_j$, identifies $\partial(\sigma_j[\Omega_j])=[E_j]$, and mass convergence
allows passage to the sum.  The same interface gives another
$L^1$ filling $V=[\Ind(L,\cdot)]$ with the same boundary.  Write
$W-V=h\,dx\wedge dy$.  Its boundary is zero, hence both distributional
partial derivatives of $h$ vanish.  Mollification makes $h*\rho_\epsilon$
constant on $\R^2$; integrability forces that constant to be zero, and an
approximate identity gives $h=0$ almost everywhere.  Therefore
\[
 \Ind(L,\cdot)=\sum_j\sigma_j\mathbf1_{\Omega_j}\quad\text{a.e.}
\]
For the cited boundary identity, given a real compactly supported test form
$\omega=A\,dx+B\,dy$, take $f=A-iB$ in Cuf\'i--Verdera and take real parts;
with their $\bar\partial$ convention,
$\operatorname{Re}(f\,dz)=\omega$ and the real part of the area integrand is
$d\omega$.  This pins the orientation and sign used above.
Tonelli, the triangle inequality, the isodiametric inequality, and (5.2)
now prove (5.1).
\end{proof}

Let
\[
 \lambda_0=\frac{x\,dy-y\,dx}{2},\qquad d\lambda_0=dx\wedge dy,
\]
and define normalized lifted primitives
\[
 F_n(0)=0,
 \qquad F_n(t)-F_n(s)=\int_{[s,t]}c_n^*\lambda_0.         \tag{5.3}
\]

\begin{hypothesis}[Generalized Green formula]\label{hyp:green}
For every closed rectifiable planar curve $L$,
\[
 \int_L\lambda_0=\int_{\R^2}\Ind(L,z)\,dz.              \tag{5.4}
\]
Here $L$ has its displayed orientation, $dx\wedge dy$ is the positive
orientation of the plane, and $\Ind$ uses the corresponding
counterclockwise-positive convention.
\status{CITED: Cuf\'i--Verdera~\cite{CV}, arXiv:1306.6832v1, main theorem, PDF p.~1.
Choose $\chi\in C_c^\infty(\R^2)$ equal to one on a neighborhood containing
both the curve trace and the support of its winding function, and apply the
cited theorem to $f(z)=\chi(z)\bar z$.  On those supports $\chi$ is constant,
so all terms involving $d\chi$ vanish.  On that neighborhood,
$\bar z\,dz=d(x^2+y^2)/2+i(x\,dy-y\,dx)$ and
$\bar\partial\bar z=1$; the exact term integrates to zero and division by
$2i$ gives (5.4).  Version, page, hypotheses, and sign independently checked}
\end{hypothesis}

\begin{proposition}[Primitive compactness]\label{prop:primitive}
Using cited Interfaces~\ref{hyp:index-current} and \ref{hyp:green}, the functions
$F_n$ have a subsequence
converging locally uniformly on $\R$.
\status{PROVED from Lemma~\ref{lem:nu},
Proposition~\ref{prop:arc}, Interfaces~\ref{hyp:index-current} and
\ref{hyp:green}, and the pinned Arzel\`a--Ascoli theorem~\cite{Rudin}}
\end{proposition}

\begin{proof}
Uniform convergence places all images in one ball $B(0,R)$.  For
$0<t-s\leq\delta\leq\tfrac12$, the lifted restriction
$c_n|_{[s,t]}$ is injective: if $s\leq r<r'\leq t$ had
$c_n(r)=c_n(r')$, injectivity of the induced map on $\R/\mathbb Z$ would give
$r'-r\in\mathbb Z$, contrary to $0<r'-r\leq\tfrac12$.  Its endpoints are therefore
distinct, so close this embedded arc by its chord.  (When $s=t$, the desired
estimate is trivial.)  Equations
(5.1)--(5.4) and direct chord parametrization yield
\[
 |F_n(t)-F_n(s)|
 \leq \frac\pi4\nu(\delta)^2+\frac R2\nu(\delta).       \tag{5.5}
\]
This is a common modulus tending to zero.  A finite cover of $[0,1]$
bounds $F_n$ uniformly there, so Arzel\`a--Ascoli supplies uniform
convergence on $[0,1]$ along a subsequence.  With $A_n=F_n(1)$,
\[
 F_n(k+r)=kA_n+F_n(r),\qquad k\in\mathbb Z,\quad0\leq r<1,
\]
because the pulled-back integrand is periodic.  This upgrades convergence
to local uniform convergence on $\R$.
\end{proof}

Pass to this subsequence and relabel it again as $(c_n,F_n)$ for the
remainder of the argument.

For the alternative Liouville form $\lambda_{\mathrm{AI}}=y\,dx$,
\[
 y\,dx=-\lambda_0+d(xy/2).                               \tag{5.6}
\]
Thus its normalized primitive is
\[
 G_n(t)=-F_n(t)
 +\frac{x_n(t)y_n(t)-x_n(0)y_n(0)}2.                     \tag{5.7}
\]
By (3.5), $c_n\to c$ uniformly, hence the coordinate products
$x_ny_n\to xy$ uniformly on $[0,1]$ (and on $\R$ by periodicity), while
$x_n(0)y_n(0)=x(0)y(0)$.  Equations (5.6)--(5.7) therefore show that the
exact-form correction converges locally uniformly; together with the
relabeled local uniform convergence of $F_n$, this proves local uniform
convergence of $G_n$.
\status{PROVED elementary exact-form conversion}

\section{Prescribed-angle rectangle theorem}

\begin{hypothesis}[Asano--Ike approximation criterion]\label{hyp:ai}
Asano--Ike~\cite{AI}, Theorem~1.1 (arXiv:2412.21057v3, PDF p.~2), uses
$e\colon\R\to\R/2\pi\mathbb Z$ for the quotient map and assumes:
\begin{enumerate}
\item a Jordan parametrization $c\colon\R/2\pi\mathbb Z\to\R^2$;
\item regular smooth Jordan parametrizations $c_n\to c$ parametrically in
  $C^0$; and
\item locally uniform convergence on $\R$ of normalized lifted primitives
  of $(c_n\circ e)^*\lambda$, for the required Liouville form $\lambda$.
\end{enumerate}
These hypotheses imply a $\theta$-rectangle for every $\theta\in(0,\pi)$:
a Euclidean rectangle whose diagonals meet at angle $\theta$.  Replacing
$\theta$ by $\pi-\theta$ gives the same unoriented diagonal angle.  Our
$\R/\mathbb Z$ parametrization is transferred to the source convention by
the linear change of parameter, which preserves both kinds of convergence.
\status{CITED; exact version, theorem, page, and hypotheses independently
checked against the primary PDF.  See
\url{https://arxiv.org/abs/2412.21057v3}}
\end{hypothesis}

\begin{hypothesis}[Explicit nondegenerate conclusion $\mathrm{AI}_{ND}$]
\label{hyp:aind}
The non-trivial intersection supplied in the source lies outside its
diagonal $\Delta_C$.  Under the source's rectangle correspondence this is a
nondegenerate $\theta$-rectangle.  At $\theta=\pi/2$ its four rectangle
vertices are pairwise distinct and its diagonals are perpendicular.
\status{CITED: definition and non-trivial/off-diagonal discussion on PDF
pp.~2--3; proof of Theorem~4.1, PDF p.~19.  Primary version:
\url{https://arxiv.org/abs/2412.21057v3}}
\end{hypothesis}

\begin{lemma}[Square criterion]
A nondegenerate Euclidean rectangle whose diagonals are perpendicular is a
positive-size square.
\status{PROVED}
\end{lemma}
\begin{proof}
For adjacent side vectors $u,v$ with $u\perp v$, the diagonal vectors are
$u+v$ and $u-v$.  Their scalar product is $|u|^2-|v|^2$; perpendicularity
therefore gives $|u|=|v|>0$.
\end{proof}

\begin{theorem}[Analytic reduction to the cited rectangle criterion]
If $c$ is a Jordan parametrization with
$V_0(c)$, then $c$ satisfies the approximation and primitive-convergence
hypotheses listed in Interface~\ref{hyp:ai}.  Consequently $c$ has
\emph{exactly the conclusion stated by the verified version} of the cited
Asano--Ike theorem.
\status{PROVED using the independently verified cited interfaces}
\end{theorem}

\begin{proof}
Interface~\ref{hyp:bg} and the proved interpolation estimate produce the
simple polygons $P_n$.  Lemma~\ref{lem:smooth} produces regular smooth
Jordan $c_n$ satisfying (3.5).  In particular $c_n\to c$ parametrically in
$C^0$, as required by Interface~\ref{hyp:ai}, and the common ball used in
Proposition~\ref{prop:primitive} follows from this uniform convergence.
Lemma~\ref{lem:nu} gives the common local
variation modulus.  Proposition~\ref{prop:primitive}, Interface
~\ref{hyp:green}, and the uniformly convergent $x_ny_n$ exact-form correction
in (5.6)--(5.7) give the locally uniformly convergent lifted primitives.
These are precisely the data listed in Interface
~\ref{hyp:ai}; therefore only the verified conclusion of that interface may
be imported.
\end{proof}

\begin{corollary}[Positive square]\label{cor:positive-square}
If $c$ is a Jordan parametrization with $V_0(c)$, then $c$ inscribes a
positive-size square with four distinct vertices.
\status{PROVED from the analytic reduction and the cited AI conclusion}
\end{corollary}
\begin{proof}
The preceding theorem verifies the hypotheses of Asano--Ike Theorem~1.1.
Interface~\ref{hyp:aind} at $\theta=\pi/2$ gives four distinct vertices of a
nondegenerate rectangle on $c$ with perpendicular diagonals.  The square
criterion makes it a positive-size square.
\end{proof}

\appendix

\section{Dependency and verification status}

Every assertion above carries one of the following audit statuses.
\begin{itemize}
\item \texttt{PROVED}: proved in this article from earlier displayed
  definitions and stated standard facts.
\item \texttt{CITED}: imported from a named external source; citation
  fidelity remains part of the verification obligation.
\item \texttt{STANDARD-CITED}: an exact, independently checked classical
  theorem is used but not reproved foundationally here.
\end{itemize}
The dependency chain is
\[
\begin{gathered}
 V_0\longrightarrow\text{polygonal approximation}
 \longrightarrow\text{regular smoothing}\\
 \longrightarrow\text{local variation}
 \longrightarrow\text{primitive compactness}
 \longrightarrow\text{Asano--Ike}.
\end{gathered}
\]
The internal arrows are proved in the article; the final status remains
conditional on the cited interfaces.  The following ledger makes the split
explicit.

\begin{center}
\begin{tabular}{>{\raggedright\arraybackslash}p{0.25\textwidth}
>{\raggedright\arraybackslash}p{0.19\textwidth}
>{\raggedright\arraybackslash}p{0.44\textwidth}}
Component & Status & Exact interface used \\
\hline
Fine-mesh/local modulus & PROVED & Definition $V_0$; seam factor $\sqrt2$ \\
PL interpolation error & PROVED & Every partition; constant $\sqrt8$ \\
Simple PL Jordan choice & CITED & BG theorem, parameter-mesh version \\
Regular smoothing & PROVED & anchored regular $C^\infty$ embedding;
$Q(0)=P(0)$ and $\|Q-P\|_\infty\leq\|Q-P\|_{1\text{-var}}$ in
(3.1)--(3.4) \\
Whole-family modulus & PROVED & tail-first then finite-prefix quantifiers \\
Arc winding estimate & PROVED & cited index/Jordan facts + internal excursion proof \\
Generalized Green & CITED & arbitrary closed rectifiable planar curve \\
Primitive compactness & PROVED & cited index/Green interfaces,
verified short lifted arcs, common modulus, Arzel\`a--Ascoli \\
Liouville conversion & PROVED & equations (5.6)--(5.7) \\
AI approximation criterion & CITED & arXiv:2412.21057v3, Thm.~1.1, p.~2 \\
Four distinct vertices & CITED & pp.~2--3 and Thm.~4.1 proof, p.~19 \\
$\theta=\pi/2$ square & PROVED & cited conclusion + perpendicular-diagonal lemma \\
\end{tabular}
\end{center}

\section{Reproducibility}

From the repository root, compile twice to resolve cross-references:
\begin{verbatim}
cd problems/square-peg/formalization
pdflatex -interaction=nonstopmode -halt-on-error \
  c02var-prescribed-angle-formalization.tex
pdflatex -interaction=nonstopmode -halt-on-error \
  c02var-prescribed-angle-formalization.tex
\end{verbatim}
The reviewable artifacts are the source
\texttt{c02var-prescribed-angle-formalization.tex}, its generated PDF, and
\texttt{ISSUE-134-AUDIT-REPAIR.md}.  The companion Lean module formalizes
only the finite and algebraic core described in \texttt{LEAN\_DAG.md}.
For each committed review checkpoint, the exact article revision and hashes
of the deliberately committed artifacts are recorded in
\texttt{ISSUE-134-AUDIT-REPAIR.md}.

\section{Limitations and governance}

This article neither proves the unrestricted Square Peg conjecture nor
promotes PR~\#122 beyond \texttt{sketch}.  It does not use PRs \#115,
\#118, or \#119 as premises.  The analytic theorem is not claimed as
\texttt{verified:lean}; the Lean module covers only identified finite and
algebraic subclaims.  The result remains conditional on the exact cited
interfaces and requires the verification-critical review prescribed by the
root \texttt{RULES.md}.  It is not a full solution of Toeplitz's conjecture.

\begin{thebibliography}{99}

\bibitem{AI}
Y.~Asano and Y.~Ike,
\emph{The rectifiable rectangular peg problem},
arXiv:2412.21057v3 (2024), Theorem~1.1, PDF p.~2; definition and
nondegeneracy discussion, pp.~2--3; proof of Theorem~4.1, p.~19.
\url{https://arxiv.org/abs/2412.21057v3}

\bibitem{Bieberbach}
L.~Bieberbach,
\emph{\"Uber eine Extremaleigenschaft des Kreises},
Jahresbericht der Deutschen Mathematiker-Vereinigung 24 (1915), 247--250.

\bibitem{BG}
H.~Boedihardjo and X.~Geng,
\emph{Simple Piecewise Geodesic Interpolation of Simple and Jordan Curves
with Applications},
arXiv:1309.1576v2 (2013), Theorem~2.2, PDF pp.~6--7.
\url{https://arxiv.org/abs/1309.1576v2}

\bibitem{CC}
J.~J. Carmona and J.~Cuf\'i,
\emph{The calculation of the $L^2$-norm of the index of a plane curve and
related formulae}, J. Analyse Math. 120 (2013), 225--253,
doi:10.1007/s11854-013-0019-9; preprint Theorem~2, PDF p.~8.

\bibitem{CV}
J.~Cuf\'i and J.~Verdera,
\emph{A generalization of Green's formula},
arXiv:1306.6832v1 (2013), main theorem, PDF p.~1.
\url{https://arxiv.org/abs/1306.6832v1}

\bibitem{GKR}
J.~Gilman, I.~Kra, and R.~E. Rodr\'iguez,
\emph{Complex Analysis: In the Spirit of Lipman Bers},
Graduate Texts in Mathematics 245, Springer, 2007,
doi:10.1007/978-0-387-74715-6, Theorem~5.23.

\bibitem{Rudin}
W.~Rudin, \emph{Principles of Mathematical Analysis}, 3rd ed.,
McGraw--Hill, 1976, Theorem~7.25, p.~158.

\end{thebibliography}

\end{document}
